%%%%%%%%%%  scan idx 330  |  folio 317  %%%%%%%%%%
\(\lambda\)-anneaux.
\item
Soit \(A\) un anneau, \(X=\operatorname{Spec}A\); soit d'autre part \(G\) un groupe, opérant trivialement sur \(A\)~; si on prend pour \((T,A)\) le topos des faisceaux d'ensembles sur \(X\), sur lesquels \(G\) opère, l'anneau \(K\) ainsi obtenu est l'anneau des représentations de \(G\) dans les \(A\)-modules projectifs de type fini.
\end{itemize}

\textbf{2.3.} Soit \(K\) un anneau commutatif à élément unité~; on se propose de munir le groupe abélien \(\Ghat(K)=1+K[[t]]^+\) d'une structure de pré-\(\lambda\)-anneau, fonctorielle par rapport à \(K\). A cet effet, nous allons appliquer les résultats de (1.10) sur le foncteur \(\Ghat\).

On définit d'abord la multiplication en se donnant un morphisme de \(\mathbb Z\)-bilinéaire de foncteurs~: \(\Ghat\times\Ghat\longrightarrow\Ghat\), par élément \(1+XYt\) de \(\Ghat(\mathbb Z[X,Y])\)~; cet élément vérifie en effet 1.7.i), et pour voir qu'il vérifie 1.7.ii), il faut montrer que dans \(\prod_{i=1}^{n}(1+T_iYt)\), exprimé au moyen des fonctions symétriques élémentaires \(X_j\) des \(T_i\), \(X_n\) ne figure pas dans les coefficients de \(t^k\) pour \(k<n\), ce qui est évident pour des raisons de degré.

La loi ainsi définie sera notée par le signe \(\circ\)~; on a donc par définition :
\begin{equation}\tag{2.3.1}
  (1+Xt)\circ(1+Yt)=1+XYt.
\end{equation}
Elle est distributive par construction, associative et commutative par les diagrammes~:
\[
\begin{tikzcd}[column sep=tiny]
\Ghat\times\Ghat\times\Ghat \arrow[r] \arrow[d] &
\Ghat\times\Ghat \arrow[d]\\
\Ghat\times\Ghat \arrow[r] & \Ghat
\end{tikzcd}
\qquad
\begin{tikzcd}[column sep=tiny]
\Ghat\times\Ghat \arrow[r] \arrow[d] &
\Ghat\times\Ghat \arrow[d]\\
\Ghat \arrow[r,"\id"] & \Ghat,
\end{tikzcd}
\]
